\documentclass[../main.tex]{subfiles}

\begin{document}
% 単独でコンパイルした場合も、全体版と章番号をそろえる。
\ifSubfilesClassLoaded{\setcounter{chapter}{0}}{}

\chapter{はじめに}

ニューラルネットワークの学習では，損失の勾配を計算してパラメータを更新する．
この仕組みを学ぶとき，最初に習うのは，勾配と逆の方向へ学習率を掛けた分だけ進む勾配降下法や，
ミニバッチで勾配を推定する確率的勾配降下法（SGD）である．
ところが，実際の学習で使う最適化手法には，AdamWやMuonといった名前が並ぶ．
これらの手法では，計算した勾配に対して，前処理を行ってから更新に使う．


実際の学習に即して最適化を理解するには，実際に使われている最適化手法，つまり，前処理を含む更新則を常に考慮にいれるのが自然だろう．
本原稿では，この立場から，前処理勾配法を中心にニューラルネットワークの最適化を概観する．


\chapterbibliography

\end{document}
